
\documentclass[11pt]{article}

\usepackage{amsmath}
\usepackage{amssymb}
\usepackage{bm}
\usepackage{xcolor}
\usepackage{algorithm}
\usepackage{algpseudocode}
\usepackage[margin=1in]{geometry}

% Additions (2026-08-19) are set in green.
\definecolor{addgreen}{HTML}{1A7A3C}

% No paragraph indentation
\setlength{\parindent}{0pt}

\title{Density-FB Derivation}
\date{}

\begin{document}

\maketitle

\section{Successor measure}

For a policy $\pi$, define the discounted successor measure over future states as
\begin{equation}
    M^\pi(s,a,ds^+) = \mathbb E_\pi \left[ \sum_{t=0}^{\infty} \gamma^t \delta_{s_{t+1}}(ds^+) \mid s_0=s,\ a_0=a \right].
    \label{eq:M-definition}
\end{equation}

Its total mass is
\begin{align}
    \int_{\mathcal S} M^\pi(s,a,ds^+)   &=  \sum_{t=0}^{\infty}\gamma^t \\
                                        &=  \frac{1}{1-\gamma}.
    \label{eq:M-mass}
\end{align}

The successor measure satisfies the Bellman equation
\begin{equation}
    M^\pi(s,a,ds^+) = P(ds^+\mid s,a) + \gamma \mathbb E_{\substack{ s'\sim P(\cdot\mid s,a)\\  a'\sim\pi(\cdot\mid s') }}  \left[  M^\pi(s',a',ds^+)   \right].
    \label{eq:M-bellman}
\end{equation}

\subsection{Density with respect to the dataset marginal}

Assume that the successor measure admits a density with respect to $\rho$:
\begin{equation}
m^\pi(s,a,s^+) = \frac{dM^\pi(s,a,\cdot)}{d\rho}(s^+).
\label{eq:m-definition}
\end{equation}

Similarly, define the one-step transition density ratio
\begin{equation}
q(s,a,s^+) = \frac{dP(\cdot\mid s,a)}{d\rho}(s^+). 
\label{eq:q-definition}
\end{equation}

Dividing Equation~\eqref{eq:M-bellman} by $\rho$ gives
\begin{equation}
m^\pi(s,a,s^+) = q(s,a,s^+) + \gamma \mathbb E_{\substack{ s'\sim P(\cdot\mid s,a)\\ a'\sim\pi(\cdot\mid s') }} \left[ m^\pi(s',a',s^+) \right].
\label{eq:m-bellman}
\end{equation}

\subsection{Squared Bellman objective}

Using a fixed target density $\bar m$, consider the squared $L_2(\rho)$ Bellman objective
\begin{align}
\mathcal L_M    =   \frac12 \mathbb E_{\substack{(s,a,s')\sim\mathcal D\\ a'\sim\pi(\cdot\mid s')\\ s^+\sim\rho }}  \left[  m(s,a,s^+)  -   q(s,a,s^+)  -   \gamma\bar m(s',a',s^+) \right]^2.
\label{eq:M-L2}
\end{align}

Expanding the square,
\begin{equation}
\begin{aligned}
\mathcal L_M    ={}&    \frac12 \mathbb E \left[ m(s,a,s^+)^2 \right]
                    -  \mathbb E \left[  m(s,a,s^+)  q(s,a,s^+)  \right] \\
                    &-  \gamma  \mathbb E \left[m(s,a,s^+) \bar m(s',a',s^+) \right]  + C,
\end{aligned}
\label{eq:M-expanded}
\end{equation}
where $C$ contains terms independent of the current density $m$.

The middle term can be rewritten using the definition of $q$:
\begin{align}
    \mathbb E_{s^+\sim\rho} \left[ m(s,a,s^+)q(s,a,s^+) \right] &= \int m(s,a,s^+) \frac{dP(\cdot\mid s,a)}{d\rho}(s^+) \,d\rho(s^+) \\
    &= \int m(s,a,s^+) \,dP(s^+\mid s,a) \\
    &= \mathbb E_{s'\sim P(\cdot\mid s,a)} \left[ m(s,a,s') \right].
    \label{eq:change-measure-M}
\end{align}

Therefore, after dropping $C$ and changing sign, the objective is
\begin{equation}
\boxed{
\begin{aligned} 
    \max_m\quad \mathcal J_M ={}& \mathbb E_{(s,a,s')\sim\mathcal D} \left[ m(s,a,s') \right]
                                + \gamma \mathbb E_{\substack{ (s,a,s')\sim\mathcal D\\ a'\sim\pi(\cdot\mid s')\\ s^+\sim\rho }} \left[ m(s,a,s^+) \bar m(s',a',s^+) \right] \\
                                &- \frac12 \mathbb E_{\substack{ (s,a,s')\sim\mathcal D\\ s^+\sim\rho }} \left[ m(s,a,s^+)^2 \right].
\end{aligned}
}
\label{eq:M-objective}
\end{equation}

\subsection{FB factorization}

FB parameterizes the successor density as
\begin{equation}
    m_z(s,a,s^+) = F(s,a,z)^\top B(s^+).
    \label{eq:fb-factorization}
\end{equation}

For the standard FB policy family, the bootstrap action is
\begin{equation}
    a'_z \in \arg\max_{a'} \bar F(s',a',z)^\top z.
\label{eq:fb-action}
\end{equation}

The target successor density is therefore
\begin{equation}
\bar m_z(s',a'_z,s^+) = \bar F(s',a'_z,z)^\top \bar B(s^+).
\end{equation}

Substituting the FB factorization into
Equation~\eqref{eq:M-objective} gives
\begin{align}
\label{eq:fb-objective}
\boxed{
\begin{aligned}
    \max_{F,B}\quad \mathcal J_{\mathrm{FB}} 
        ={}& \mathbb E_{\substack{ z\sim p(z)\\ (s,a,s')\sim\mathcal D }} \left[ F(s,a,z)^\top B(s') \right] 
        + \gamma \mathbb E_{\substack{ z\sim p(z)\\ (s,a,s')\sim\mathcal D\\ s^+\sim\rho}} \left[ \left( F(s,a,z)^\top B(s^+) \right) \left( \bar F(s',a'_z,z)^\top \bar B(s^+) \right) \right] \\
        &- \frac12 \mathbb E_{\substack{ z\sim p(z)\\ (s,a,s')\sim\mathcal D\\ s^+\sim\rho}} \left[ \left( F(s,a,z)^\top B(s^+) \right)^2 \right].
\end{aligned}
}
\end{align}

\section{Density-FB}

\subsection*{Step 1: normalize the successor measure}

Define the normalized successor measure
\begin{equation}
\boxed{ \nu^\pi(s,a,ds^+) = (1-\gamma) M^\pi(s,a,ds^+). }
\label{eq:nu-definition}
\end{equation}

Therefore $\nu^\pi$ is a probability measure over successor states.
Multiply Equation~\eqref{eq:M-bellman} by $1-\gamma$:
\begin{equation}
    (1-\gamma)M^\pi(s,a,ds^+) = (1-\gamma)P(ds^+\mid s,a)
                                + \gamma \mathbb E_{\substack{ s'\sim P(\cdot\mid s,a)\\  a'\sim\pi(\cdot\mid s') }} \left[ (1-\gamma) M^\pi(s',a',ds^+) \right].
\end{equation}

Using Equation~\eqref{eq:nu-definition},
\begin{equation}
\boxed{ 
\nu^\pi(s,a,ds^+) = (1-\gamma)P(ds^+\mid s,a) 
                        + \gamma \mathbb E_{\substack{ s'\sim P(\cdot\mid s,a)\\ a'\sim\pi(\cdot\mid s') }}
                                \left[ \nu^\pi(s',a',ds^+) \right].
}
\label{eq:nu-bellman}
\end{equation}

The immediate component has mass $1-\gamma$ and the bootstrap component has mass $\gamma$.

\subsection*{Step 2: density of the normalized measure}

Define
\begin{equation}
    h^\pi(s,a,s^+) = \frac{d\nu^\pi(s,a,\cdot)}{d\rho}(s^+). 
    \label{eq:h-definition}
\end{equation}

Because $\nu^\pi$ is normalized,
\begin{equation}
\boxed{ \mathbb E_{s^+\sim\rho} \left[ h^\pi(s,a,s^+) \right] = 1 }
\label{eq:h-normalized}
\end{equation}

Using Equation~\eqref{eq:q-definition}, the normalized Bellman equation becomes
\begin{equation}
h^\pi(s,a,s^+) = (1-\gamma)q(s,a,s^+) + \gamma \mathbb E_{\substack{ s'\sim P(\cdot\mid s,a)\\ a'\sim\pi(\cdot\mid s') }}
                                            \left[ h^\pi(s',a',s^+) \right]. \label{eq:h-bellman}
\end{equation}

For a sampled transition and a target-policy bootstrap action $a'\sim\bar\pi(\cdot\mid s'),$ define the target density (its dependence on the sampled $(s',a')$ is left implicit in the notation)
\begin{equation}
    \boxed{ h^{\mathrm{tgt}}(s,a,s^+) = (1-\gamma)q(s,a,s^+) + \gamma \bar h(s',a',s^+). }
    \label{eq:h-target}
\end{equation}

The target is normalized:
\begin{equation}
\begin{aligned}
\mathbb E_{s^+\sim\rho} \left[ h^{\mathrm{tgt}}(s,a,s^+) \right] &= (1-\gamma) \mathbb E_{s^+\sim\rho} [q(s,a,s^+)]
                                                                    + \gamma \mathbb E_{s^+\sim\rho} [\bar h(s',a',s^+)] \\
                                                                 &= (1-\gamma) + \gamma = 1.
\end{aligned}
\end{equation}

\subsection*{Step 3: fit the normalized Bellman target}

At this point both $h$ and $h^{\mathrm{tgt}}$ are normalized densities. 
We fit the target distribution with cross-entropy by maximizing
\begin{equation}
    \mathcal J_{\mathrm{DFB}} = \mathbb E_{s^+\sim\rho} \left[ h^{\mathrm{tgt}}(s,a,s^+) \log h(s,a,s^+) \right].
\label{eq:dfb-ce}
\end{equation}

Substituting Equation~\eqref{eq:h-target},
\begin{equation}
\begin{aligned}
    \mathcal J_{\mathrm{DFB}} ={}& (1-\gamma) \mathbb E_{s^+\sim\rho} \left[ q(s,a,s^+) \log h(s,a,s^+) \right] \\
                                &+ \gamma \mathbb E_{s^+\sim\rho} \left[ \bar h(s',a',s^+) \log h(s,a,s^+) \right].
\end{aligned}
    \label{eq:dfb-expanded}
\end{equation}

The immediate term can again be rewritten by change of measure:
\begin{align}
    \mathbb E_{s^+\sim\rho} \left[ q(s,a,s^+) \log h(s,a,s^+) \right] &= \int \frac{dP(\cdot\mid s,a)}{d\rho}(s^+) \log h(s,a,s^+) \,d\rho(s^+) \\
    &= \int \log h(s,a,s^+) \,dP(s^+\mid s,a) \\
    &= \mathbb E_{s'\sim P(\cdot\mid s,a)} \left[ \log h(s,a,s')\right].
\label{eq:change-measure-h}
\end{align}

Hence the Density-FB objective is
\begin{equation}
\label{eq:dfb-population}
\boxed{
\begin{aligned}
    \max_h\quad \mathcal J_{\mathrm{DFB}} ={}& (1-\gamma) \mathbb E_{(s,a,s')\sim\mathcal D} \left[ \log h(s,a,s') \right] \\
                                            &+ \gamma \mathbb E_{\substack{ (s,a,s')\sim\mathcal D\\ a'\sim\bar\pi(\cdot\mid s')\\ s^+\sim\rho }}
                                                \left[ \bar h(s',a',s^+) \log h(s,a,s^+) \right].
\end{aligned}
}
\end{equation}

{\color{addgreen}
\subsection*{Soft immediate targets (the flagship variant)}

The target of Equation~\eqref{eq:h-target} places the immediate mass $1-\gamma$ exactly on the observed successor.
The soft variant smooths it over a neighborhood.
Let $k_\epsilon(x,u)>0$ be a fixed kernel and define the base-measure-dependent Markov kernel
\begin{equation}
    K_{\rho,\epsilon}(du\mid x)
    =
    \frac{k_\epsilon(x,u)\,\rho(du)}{\int k_\epsilon(x,v)\,\rho(dv)}.
\label{eq:soft-kernel}
\end{equation}
Replacing the immediate term of Equation~\eqref{eq:nu-bellman} by
$\mathbb E_{s'\sim P(\cdot\mid s,a)}\left[K_{\rho,\epsilon}(\cdot\mid s')\right]$
leaves a $\gamma$-contraction whose fixed point is the kernel-smoothed occupancy
\begin{equation}
    \nu_K^\pi(du\mid s,a)
    =
    (1-\gamma)\sum_{t=0}^{\infty}\gamma^t\,
    \mathbb E\left[ K_{\rho,\epsilon}(du\mid s_{t+1}) \mid s_0=s,\ a_0=a,\ \pi \right].
\label{eq:soft-fixed-point}
\end{equation}
For any reward $r$,
\begin{equation}
    \int r(u)\, \nu_K^\pi(du\mid s,a)
    =
    \int (K_{\rho,\epsilon} r)(x)\, \nu^\pi(dx\mid s,a),
    \qquad
    (K_{\rho,\epsilon} r)(x) = \int r(u)\, K_{\rho,\epsilon}(du\mid x),
\label{eq:soft-smoothed-reward}
\end{equation}
so the soft model evaluates the original occupancy under the \emph{smoothed} reward $K_{\rho,\epsilon}r$, not under $r$ itself.
Sums of rewards remain linear under smoothing, but sharp boundaries, maxima, and exclusion regions are softened; the hard target is the semantically exact interface for such rewards.
Note also that $K_{\rho,\epsilon}$ is the kernel multiplied by the data density, not a Lebesgue kernel: its semantics depend on $\rho$.
The flagship experiments run this soft variant; the hard target is the ablation.
}

\subsection*{Step 4: parameterize the normalized density}

The density must satisfy $h(s,a,s^+)\geq 0$ and $\mathbb E_{s^+\sim\rho} [h(s,a,s^+)] = 1.$
Define the bilinear score
\begin{equation}
f_\theta(s,a,z,s^+) = \frac{ F_\theta(s,a,z)^\top B_\theta(s^+) }{\tau}.
\label{eq:score}
\end{equation}

A natural parameterization of this is
\begin{equation}
    \boxed{ h_\theta(s,a,z,s^+) = \frac{ \exp f_\theta(s,a,z,s^+) }{ Z_\theta(s,a,z) }, }
\label{eq:h-exponential}
\end{equation}
where
\begin{equation}
Z_\theta(s,a,z) = \mathbb E_{s^+\sim\rho} \left[ \exp f_\theta(s,a,z,s^+) \right].
\label{eq:partition}
\end{equation}

The corresponding target density is constructed from target networks:
\begin{equation}
    \boxed{ \bar h(s',a',z,s^+) = \frac{ \exp \bar f(s',a',z,s^+) }{ \bar Z(s',a',z) }, }
\end{equation}
where
\begin{equation}
    \bar f(s',a',z,s^+) = \frac{ \bar F(s',a',z)^\top \bar B(s^+) }{\tau},
\end{equation}
and
\begin{equation}
    \bar Z(s',a',z) = \mathbb E_{s^+\sim\rho} \left[ \exp \bar f(s',a',z,s^+) \right].
\end{equation}

The normalization of the current density follows directly:
\begin{equation}
\mathbb E_{s^+\sim\rho} [h_\theta(s,a,z,s^+)]   = \frac{ \mathbb E_{s^+\sim\rho} [\exp f_\theta(s,a,z,s^+)] }{ Z_\theta(s,a,z) }
                                                = 1.
\end{equation}

Thus the role of the FB bilinear quantity changes.
In ordinary FB, $F(s,a,z)^\top B(s^+)$ directly represents the successor density.
In Density-FB, $\frac{F(s,a,z)^\top B(s^+)}{\tau}$ is a log-score whose exponential is normalized with respect to $\rho$.

\subsection*{Step 5: substitute the exponential density}

From Equation~\eqref{eq:h-exponential},
\begin{equation}
    \log h_\theta(s,a,z,s^+) = f_\theta(s,a,z,s^+) - \log Z_\theta(s,a,z).
    \label{eq:log-h}
\end{equation}

Substitute this into Equation~\eqref{eq:dfb-population}:
\begin{align}
    \mathcal J_{\mathrm{DFB}} ={}& (1-\gamma) \mathbb E_{\mathcal D} \left[ f_\theta(s,a,z,s') - \log Z_\theta(s,a,z) \right] \nonumber\\
                                &+ \gamma \mathbb E_{\substack{ (s,a,s')\sim\mathcal D\\ a'\sim\bar\pi(\cdot\mid s',z)\\ s^+\sim\rho }}
                                        \left[ \bar h(s',a',z,s^+) \left( f_\theta(s,a,z,s^+) - \log Z_\theta(s,a,z) \right) \right].
    \label{eq:dfb-substitution}
\end{align}

Expanding,
\begin{align}
    \mathcal J_{\mathrm{DFB}} ={}& (1-\gamma) \mathbb E_{\mathcal D} [f_\theta(s,a,z,s')]
                                - (1-\gamma) \mathbb E_{\mathcal D} [\log Z_\theta(s,a,z)] \nonumber\\
                                &+ \gamma \mathbb E_{\substack{ (s,a,s')\sim\mathcal D\\ a'\sim\bar\pi(\cdot\mid s',z)\\ s^+\sim\rho }}
                                        \left[ \bar h(s',a',z,s^+) f_\theta(s,a,z,s^+) \right]
                                - \gamma \mathbb E_{\substack{ (s,a,s')\sim\mathcal D\\ a'\sim\bar\pi(\cdot\mid s',z)\\ s^+\sim\rho }}
                                        \left[ \bar h(s',a',z,s^+) \log Z_\theta(s,a,z) \right].
    \label{eq:dfb-four-terms}
\end{align}

Since $\log Z_\theta(s,a,z)$ does not depend on $s^+$,
\begin{equation}
    \mathbb E_{s^+\sim\rho} \left[ \bar h(s',a',z,s^+) \log Z_\theta(s,a,z) \right] = \log Z_\theta(s,a,z) \mathbb E_{s^+\sim\rho} \left[ \bar h(s',a',z,s^+) \right]
    = \log Z_\theta(s,a,z).
\end{equation}

Therefore the two normalization terms combine as
\begin{equation}
-(1-\gamma)\log Z - \gamma\log Z = -\log Z.
\end{equation}

The final population objective is
\begin{align}
\label{eq:dfb-final}
\boxed{
\begin{aligned}
    \max_\theta\quad \mathcal J_{\mathrm{DFB}} ={}& (1-\gamma) \mathbb E_{\mathcal D} \left[ f_\theta(s,a,z,s') \right] 
                                                + \gamma \mathbb E_{\substack{(s,a,s')\sim\mathcal D\\a'\sim\bar\pi(\cdot\mid s',z)\\ s^+\sim\rho}} \left[ \bar h(s',a',z,s^+)f_\theta(s,a,z,s^+)\right] \\
                                                &- \mathbb E_{\mathcal D} \left[\log\mathbb E_{s^+\sim\rho}\exp f_\theta(s,a,z,s^+)\right].
\end{aligned}
}
\end{align}

Or if we are to fully expand this objective with the $F,B$ embeddings, we get

\begin{align}
\boxed{
\begin{aligned}
    \max_{F,B}\quad \mathcal J_{\mathrm{DFB}} ={}& (1-\gamma) \mathbb E_{\substack{ z\sim p_Z\\ (s,a,s')\sim\mathcal D }} \left[                                                               \frac{F(s,a,z)^\top B(s')}{\tau} \right] \\
                                                &+ \gamma \mathbb E_{\substack{z\sim p_Z\\(s,a,s')\sim\mathcal D\\a'\sim\bar\pi(\cdot\mid s',z)\\s^+\sim\rho}} \left[ \frac{\exp\left(\frac{\bar F(s',a',z)^\top\bar B(s^+)}{\tau}\right)}{\mathbb E_{\tilde s^+\sim\rho}\left[\exp\left(\frac{\bar F(s',a',z)^\top\bar B(\tilde s^+)}{\tau}\right) \right]} \frac{F(s,a,z)^\top B(s^+)}{\tau}\right] \\
                                                &- \mathbb E_{\substack{z\sim p_Z\\(s,a,s')\sim\mathcal D}}\left[\log\mathbb E_{s^+\sim\rho}\exp\left(\frac{F(s,a,z)^\top B(s^+)}{\tau}\right)\right].
\end{aligned}
}
\label{eq:dfb-final-expanded}
\end{align}

\subsection{Importance-weight connection}

The normalized density
\begin{equation}
    h(s,a,z,s^+) = \frac{d\nu^\pi(s,a,\cdot)}{d\rho}(s^+)
\end{equation}
is also an importance ratio between the normalized successor distribution $\nu^\pi$ and the reference distribution $\rho$.
For any function $g$,
\begin{align}
    \mathbb E_{s^+\sim\nu^\pi(\cdot\mid s,a,z)} \left[g(s^+)\right] &= \int g(s^+) h(s,a,z,s^+) \, d\rho(s^+) \\
                                                                    &= \mathbb E_{s^+\sim\rho} \left[ h(s,a,z,s^+)g(s^+) \right].
\end{align}

Under the exponential parameterization, define the unnormalized importance weight
\begin{equation}
    w(s,a,z,s^+) = \exp\left( \frac{F(s,a,z)^\top B(s^+)}{\tau} \right).
    \label{eq:importance-weight}
\end{equation}

Then
\begin{equation}
    \boxed{
    h(s,a,z,s^+) = \frac{w(s,a,z,s^+)}{ \mathbb E_{\tilde s^+\sim\rho} [w(s,a,z,\tilde s^+)] }. 
    }
    \label{eq:importance-ratio}
\end{equation}

Similarly, the target networks define
\begin{equation}
    \bar w(s',a',z,s^+) = \exp\left( \frac{ \bar F(s',a',z)^\top\bar B(s^+)}{\tau}\right),
\end{equation}
and
\begin{equation}
    \boxed{
    \bar h(s',a',z,s^+) = \frac{ \bar w(s',a',z,s^+)}{\mathbb E_{\tilde s^+\sim\rho}[\bar w(s',a',z,\tilde s^+)]}.
    }
    \label{eq:target-importance-ratio}
\end{equation}

Therefore the Density-FB objective can be written directly as an importance-weighted objective:
\begin{align}
\boxed{
\begin{aligned}
    \max_{F,B}\quad \mathcal J_{\mathrm{DFB}} ={}& (1-\gamma) \mathbb E_{\substack{z\sim p_Z\\ (s,a,s')\sim\mathcal D}} 
                                                    \left[ \log \frac{w(s,a,z,s')}{\mathbb E_{s^+\sim\rho}[w(s,a,z,s^+)]}\right] \\
                                                &+ \gamma \mathbb E_{\substack{z\sim p_Z\\(s,a,s')\sim\mathcal D\\a'\sim\bar\pi(\cdot\mid s',z)\\s^+\sim\rho}} \left[ \frac{\bar w(s',a',z,s^+)}{\mathbb E_{\tilde s^+\sim\rho}[\bar w(s',a',z,\tilde s^+)]} \log\frac{w(s,a,z,s^+)}{\mathbb E_{\tilde s^+\sim\rho}[w(s,a,z,\tilde s^+)]}\right].
\end{aligned}
}
\label{eq:dfb-importance-population}
\end{align}

The first term evaluates the likelihood of the observed successor $s'$.
In the second term, states are sampled from $\rho$ and reweighted by the target importance ratio
\[
\bar h(s',a',z,s^+) = \frac{d\bar\nu(s',a',z,\cdot)}{d\rho}(s^+),
\]
so that the expectation corresponds to an expectation under the bootstrapped successor distribution.

{\color{addgreen}
When the reference expectations here are replaced by finite pools of samples from $\rho$, both ratios become \emph{self-normalized importance-sampling} (SNIS) estimators: consistent as the pool grows, biased at any finite pool size, with the effective sample size of the weights as the standard collapse diagnostic.

\subsection{Relation to TD-InfoNCE}

The update derived above is closely related to TD-InfoNCE (Zheng, Salakhutdinov, and Eysenbach, ICLR 2024; arXiv:2310.20141).
With hard immediate labels and matched query supports, the cross-entropy update of Equation~\eqref{eq:dfb-population} reduces to the categorical Bellman update underlying TD-InfoNCE: a $(1-\gamma)$-weighted classification of the observed next state against samples from the marginal, plus a $\gamma$-weighted classification term soft-labeled by the stop-gradient target softmax.
Their analysis likewise identifies the induced classifier with the normalized discounted future-state distribution --- the object $h$ of Step~2.
We therefore do not claim the categorical update itself as a contribution.

The differences are in what surrounds the update.
TD-InfoNCE has no task latent: its critic is goal-conditioned, $f(s,a,g,s^+)$, so there is no $F(s,a,z)^\top B(s^+)$ policy family and no mechanism for inferring a policy index from a reward; its arbitrary-reward evaluation is a fixed-policy observation, not a deployed interface.
Its actor maximizes the classification likelihood of the commanded goal, rather than integrating externally supplied reward samples (Section~\ref{sec:policy-extraction}).
It has no kernel-smoothed variant (the soft immediate target above).
TD-InfoNCE does operate offline --- D4RL AntMaze with a behavioral-cloning regularizer --- so the distinction is the reward-prompted task interface, not the offline setting.

The distinction is also empirical: a paper-exact TD-InfoNCE control (their normalized towers, query supports, goal-contrastive actor, and behavioral-cloning regularizer, on the same outcome coordinates) reaches $0.576$ on cube-double, where soft Density-FB reaches $0.736$ --- one seed each, preliminary.

\section{Turning Density-FB into an algorithm}
\label{sec:algorithm}

Consider a minibatch $\{(s_i,a_i,s_i')\}_{i=1}^{n}$.
The batch next states $\{s_j'\}_{j=1}^{n}$ are used as the finite candidate successor-state pool.
For every transition $i$ and candidate successor state $s_j'$, compute
\begin{equation}
    f_{ij} = \frac{ F_\theta(s_i,a_i,z_i)^\top B_\theta(s_j') }{\tau},
\end{equation}
and normalize over the candidate states:
\begin{equation}
\boxed{
    p_{ij} = \frac{ \exp(f_{ij}) }{ \sum_{k=1}^{n}\exp(f_{ik}) }.
}
\label{eq:p-current}
\end{equation}

For each observed next state $s_i'$, obtain a bootstrap action $a_i' \sim \bar\pi(\cdot\mid s_i',z_i)$.
Using the target networks, compute $\bar f_{ij}' = \bar F(s_i',a_i',z_i)^\top \bar B(s_j') / \tau$ and normalize:
\begin{equation}
\boxed{
    \bar p_{ij}' = \frac{ \exp(\bar f_{ij}') }{ \sum_{k=1}^{n} \exp(\bar f_{ik}') }.
}
\label{eq:p-target-bootstrap}
\end{equation}

For transition $i$, the observed next state is $s_i'$.
Because that state appears as candidate $i$ in the query pool, define $e_{ij} = \mathbf 1[j=i]$.
The normalized Bellman target is
\begin{equation}
\boxed{
    p_{ij}^{\mathrm{tgt}} = (1-\gamma)e_{ij} + \gamma\bar p_{ij}'.
}
\label{eq:p-bellman-target}
\end{equation}
Both $e_i$ and $\bar p_i'$ are probability vectors, hence $\sum_j p_{ij}^{\mathrm{tgt}} = 1$.

The successor model is trained with cross-entropy:
\begin{equation}
\boxed{
    \mathcal L_{\mathrm{DFB}}
    =
    -\frac{1}{n} \sum_{i=1}^{n} \sum_{j=1}^{n} p_{ij}^{\mathrm{tgt}} \log p_{ij}.
}
\label{eq:dfb-algorithm-loss}
\end{equation}

A representation-learning update therefore consists of:
\begin{enumerate}
    \item Sample a minibatch $(s_i,a_i,s_i')$.
    \item Use the batch next states $s_j'$ as candidate successor states.
    \item Compute the current normalized distribution $p_{ij}$.
    \item From each $s_i'$, compute the target-network bootstrap distribution $\bar p_{ij}'$.
    \item Form $p_i^{\mathrm{tgt}} = (1-\gamma)e_i + \gamma\bar p_i'$.
    \item Minimize the cross-entropy between $p_i^{\mathrm{tgt}}$ and $p_i$.
    \item Update the target networks.
\end{enumerate}

\subsection*{Implementation notes}

Five properties of the implementation deserve explicit note.
(i) The candidate states and the backward inputs are the goal-relevant object coordinates $\phi(s')$, not raw states: the derivation applies verbatim to the pushforward measure, with $\rho$ the object-coordinate next-state marginal.
(ii) The batch softmax replaces the reference expectations of Step~4 with finite-pool estimates: the resulting objective is a self-normalized importance-sampling estimate, consistent as $n$ grows and biased at finite $n$; the effective sample size of the weights is the monitored diagnostic.
(iii) Two score heads are trained, each bootstrapping its own target head, with the mean cross-entropy; the pessimistic minimum over heads is applied only in the actor.
(iv) The bootstrap expectation over $a'$ is estimated with a single stop-gradient action drawn from the current actor.
(v) The one-hot label $e_{ij}=\mathbf 1[j=i]$ presumes distinct candidates; in dwell-heavy data, near-duplicate candidates make the label unattainable at identical logits --- an irreducible cross-entropy floor that is instrumented rather than corrected.
}

\section{Policy extraction}
\label{sec:policy-extraction}

\subsection{Standard FB}

In standard FB, the value of a reward $r$ can be obtained directly from the successor density $Q_r^{\pi_z}(s,a)$.
Define the reward task vector
\begin{equation}
    z_r = \mathbb E_{s^+\sim\rho} \left[ r(s^+)B(s^+) \right].
\label{eq:fb-zr}
\end{equation}

Setting the policy index to $z=z_r$ therefore gives
\begin{equation}
    Q_r^{\pi_{z_r}}(s,a) = F(s,a,z_r)^\top z_r.
\end{equation}

The actor can consequently be co-trained with the FB representation by maximizing
\begin{equation}
    \max_\pi \quad \mathbb E_{s,z}\left[ F(s,\pi(s,z),z)^\top z \right].
\end{equation}

\subsection{Density-FB}

In Density-FB, the normalized successor density is instead $h(s,a,z,s^+)$.
Since $\nu^\pi=(1-\gamma)M^\pi$,
\begin{align}
    (1-\gamma)Q_r^{\pi_z}(s,a)  &= \int r(s^+) \nu^{\pi_z}(s,a,ds^+) \\
                                &= \mathbb E_{s^+\sim\rho} \left[ h(s,a,z,s^+) r(s^+) \right].
\end{align}

For a task $z$ with reward $r_z$, define
\begin{equation}
\boxed{
    q_z(s,a) = \mathbb E_{s^+\sim\rho} \left[ h(s,a,z,s^+)r_z(s^+) \right].
}
\label{eq:actor-q}
\end{equation}

Thus
\begin{equation}
    q_z(s,a) = (1-\gamma)Q_{r_z}^{\pi_z}(s,a),
\end{equation}
and the factor $1-\gamma$ does not affect the greedy action.

Using the exponential parameterization of the normalized density,
\begin{equation}
    h(s,a,z,s^+)    = \frac{ \exp f(s,a,z,s^+) }{ \mathbb E_{\tilde s^+\sim\rho} \left[\exp f(s,a,z,\tilde s^+)\right]},
\end{equation}
Equation~\eqref{eq:actor-q} becomes
\begin{equation}
\boxed{
    q_z(s,a) = \frac{\mathbb E_{s^+\sim\rho}\left[\exp f(s,a,z,s^+)r_z(s^+)\right]}{\mathbb E_{s^+\sim\rho}\left[\exp f(s,a,z,s^+)\right]}.
}
\label{eq:actor-q-density}
\end{equation}

The actor now maximizes the reward integrated under the successor distribution predicted by Density-FB.

{\color{addgreen}
Two consequences follow.
First, Equation~\eqref{eq:fb-action} is no longer a policy-improvement rule: $q_z$ is nonlinear in the score $f$ through the partition function, so no closed-form greedy action in the latent survives, and the actor must maximize Equation~\eqref{eq:actor-q-density} directly.
Second, with a finite candidate pool $\{s_j^+\}_{j=1}^{M}$, $s_j^+\sim\rho$, Equation~\eqref{eq:actor-q-density} becomes the pool-softmax form
\begin{equation}
    q_z(s,a) = \sum_{j=1}^{M} p_j(s,a,z)\, r_z(s_j^+),
    \qquad
    p_j = \frac{\exp f(s,a,z,s_j^+)}{\sum_{k}\exp f(s,a,z,s_k^+)},
\label{eq:actor-q-pool}
\end{equation}
the self-normalized importance-sampling estimate of the reward under the predicted successor distribution.
The pool is the same successor support the density loss trains against, so the actor queries the model exactly where it was fit.
With twin score heads, the actor uses the elementwise minimum of the two head values; the representation loss itself averages the heads.

\subsection*{Behavior-regularized maximization}

Offline, an unconstrained $\arg\max_a q_z(s,a)$ leaves the data support.
The actor therefore maximizes Equation~\eqref{eq:actor-q} under a behavior-cloning constraint:
\begin{equation}
    \mathcal L_{\mathrm{actor}}
    =
    -\,\mathbb E_{s,z}\left[ q_z(s, \pi(s,z)) \right]
    \;+\;
    \beta\, \mathcal L_{\mathrm{BC}}(\pi),
\label{eq:actor-raw}
\end{equation}
where $\mathcal L_{\mathrm{BC}}$ is a flow-matching distillation loss (or a squared-error term for a deterministic actor) toward the dataset actions.
Throughout, the actor receives no gradient path into $F$ or $B$ (the representation is frozen with respect to the actor's update), and $r_z$ is treated as a constant.

\subsection*{The scale correction}

Equation~\eqref{eq:actor-raw} is not usable as written, because its two terms are not commensurate.
The critic $q_z$ is a probability-scale quantity: its magnitude is set by the pool normalization and by an affine offset the actor cannot influence.
Writing $q_z = c + \tilde q_z$ with $c$ constant in $a$, only $\tilde q_z$ carries policy information, and in variants where the value contains an explicit immediate term the controllable part is $(1-\gamma)$-scaled while $|q_z| = O(1)$.
Normalizing the greedy term by the raw magnitude $\mathbb E|q_z|$ --- the standard offline recipe --- then divides the useful gradient by the offset and suppresses it against $\mathcal L_{\mathrm{BC}}$.

Density-FB instead normalizes by the \emph{centered} magnitude, held constant through the gradient:
\begin{equation}
\boxed{
    \mathcal L_{\mathrm{actor}}
    =
    -\,
    \frac{
        \mathbb E\left[ q_z(s,\pi(s,z)) \right]
    }{
        \mathrm{sg}\!\left(
        \mathbb E\bigl| q_z - \mathbb E[q_z] \bigr| + \varepsilon
        \right)
    }
    \;+\;
    \beta\, \mathcal L_{\mathrm{BC}}(\pi).
}
\label{eq:actor-centered}
\end{equation}
This has two properties.
First, the greedy/BC balance is invariant to an affine offset $q_z \mapsto q_z + c$: the offset changes neither the normalizer nor the gradient of the numerator, so the effective strength of the greedy term is set by the variation of $q_z$ the actor can actually control.
Second, the centering lives entirely inside the stop-gradient, and this placement is load-bearing: if the mean is instead subtracted on the differentiated path, the mean-gradient term cancels the greedy gradient \emph{exactly} whenever the minimum over twin heads selects the same head on every sample, and attenuates it otherwise.
The correction must change the normalization of the objective, not the objective itself.

\subsection*{Summary of the modification}

Relative to the FB greedy actor, Density-FB changes three things:
(1) the maximized quantity is the reward integral of Equation~\eqref{eq:actor-q} under the normalized successor model, with the reward supplied externally, rather than the bilinear latent value $F^\top z$;
(2) maximization is behavior-regularized, with the representation frozen with respect to the actor's gradients;
(3) the greedy term is normalized by the stop-gradded centered magnitude of the critic, which makes the greedy/BC balance invariant to the probability-scale offset the normalized density introduces.
The third item is the modification that makes the greedy term effective in practice; without it, the actor collapses to behavior cloning.

\section{The task interface: where $r_z$ comes from}

\subsection*{What standard FB promises, and why it fails here}

Standard FB supports posterior reward specification: after unsupervised training, Equation~\eqref{eq:fb-zr} maps a reward to a task vector and $\pi_{z_r}$ is deployed with no further optimization.
This rests on two properties that hold jointly in FB: the value is \emph{linear} in the task vector, $Q_r^{\pi_z} = F(s,a,z)^\top \mathbb E_\rho[rB]$, and the training loss constrains $B$ itself, so Equation~\eqref{eq:fb-zr} is evaluated in meaningful coordinates.

Density-FB breaks both properties.
First, the value $(1-\gamma) Q_r = \mathbb E_{s^+\sim\rho}[h\, r]$ is nonlinear in the score through the partition function, so no closed-form inverse from rewards to task vectors follows from the model.
Second, the cross-entropy objective constrains only the \emph{normalized} density: any transformation of $B$ invisible to the softmax --- in particular the additive gauge $B \mapsto B + b$, since $F^\top b$ is constant over the candidate axis --- leaves every prediction unchanged.
The coordinates in which Equation~\eqref{eq:fb-zr} is computed are therefore arbitrary.
Empirically this is not a technicality: the uncentered second moment of the learned $B$ is dominated by a gauge direction $33$--$103\times$ larger than the informative signal, the moment embedding assigns rewards to task vectors at chance, and a constant reward shift $r \mapsto r + \beta$ --- which leaves the optimal policy unchanged --- rotates Equation~\eqref{eq:fb-zr} to a nearly orthogonal vector.

\subsection*{Three candidate interfaces}

The actor objective requires an association $z \leftrightarrow r_z$ during actor training.
There are three ways to supply it.

\textbf{(i) Goal conditioning.}
Set $z = g$ and let $r_g(\cdot) = \kappa(g,\cdot)$ be a fixed kernel around the commanded goal.
The reward is known at training time.
This is the primary instantiation in our experiments; it is goal-conditioned RL, and we claim it as such --- not as posterior reward specification.

\textbf{(ii) The linear self-supervised family.}
Set $r_z(s^+) = B(s^+)^\top z$, giving the compact actor value $q_z(s,a) = \left(\mathbb E_\rho[h(s,a,z,s^+) B(s^+)]\right)^\top z$.
This inherits FB's algebra but none of its justification: the reward family is expressed in the gauge-arbitrary coordinates above.
We tested it, and it fails: reward reconstruction from $B^\top z$ is at or below chance for uncentered task vectors, and an actor trained on this family reaches $0.368$ where the same architecture with externally supplied rewards reaches $0.812$.

\textbf{(iii) A reward-sample interface.}
Supply the reward, after training, as evaluations $\{(s_j^+, r_j)\}_{j=1}^{M}$ on states $s_j^+ \sim \rho$ --- the same posterior interface FB uses to estimate Equation~\eqref{eq:fb-zr} in practice --- and construct the task vector by \emph{regression} rather than by the moment:
\begin{equation}
\boxed{
    z_r
    =
    \sqrt d\,\operatorname{normalize}\!\left(
    \left( C_B + \lambda I \right)^{-1}
    \frac1M \sum_{j=1}^{M}
    \left( B(s_j^+) - \mu_B \right)\left( r_j - \bar r \right)
    \right),
}
\label{eq:centered-dual}
\end{equation}
with $\mu_B, C_B$ the mean and covariance of $B$ under $\rho$ and $\bar r$ the sample mean reward.
This is the ridge coefficient for reconstructing the centered reward from the centered features.
It is exactly invariant to the additive gauge $B \mapsto B+b$ and to positive affine reward transformations $r \mapsto \alpha r + \beta$ --- the invariances a task code must have and that Equation~\eqref{eq:fb-zr} provably lacks here.
The actor is trained across a family of such tasks (hindsight goal kernels and sampled task vectors), integrating the \emph{actual} reward samples in its value; the task vector is only a policy index, and reward semantics never pass through $B^\top z$.

Two structural remarks.
Because Equation~\eqref{eq:centered-dual} centers the rewards, it has no one-sample limit ($M=1$ gives $r_1 - \bar r = 0$): point-goal prompting is a separate analytic mode, $z_g = \sqrt d\,\operatorname{normalize}(B(g) - \mu_B)$, not a special case of the sample encoder.
And constant rewards give $z_r = 0$ before normalization; they are mapped to a designated neutral index, since every policy is optimal for them.

\subsection*{What is claimed}

Interface (iii) restores the \emph{form} of posterior reward specification: reward in, task vector out, frozen policy deployed.
It does not restore FB's theorem.
The deployment-time guarantee is empirical and factors into three measurable conditions: the reward must be representable in the centered feature basis (measured by held-out reconstruction), the encoder must assign it to the right index (measured against a search over task vectors using the model's own reward integral), and the trained policy family must contain a policy for it (the residual when reconstruction and assignment both succeed).
We therefore state the claim as: a gauge-corrected, affine-invariant reward-to-index interface whose zero-shot validity is established per reward family by evaluation, not by construction.

\section{The algorithm}

Algorithm~\ref{alg:dfb} assembles the pieces: the categorical representation update of Section~\ref{sec:algorithm}, the behavior-regularized actor with the centered scale correction, the target-network update, and the two deployment-time prompting modes.

\begin{algorithm}[h]
\caption{Density-FB (hard or soft immediate target)}
\label{alg:dfb}
\begin{algorithmic}[1]
\Require dataset $\mathcal D$, towers $F_\theta, B_\theta$ (two score heads), targets $\bar F, \bar B$, actor $\pi$, kernel $k_\epsilon$, temperature $\tau$, rates $\eta, \zeta$
\For{each training step}
    \State sample minibatch $\{(s_i,a_i,s_i')\}_{i=1}^{n}$; task $z_i$ (commanded goal or sampled task vector)
    \State query pool $\{u_j\}_{j=1}^{n} \gets \phi(s_j')$ \Comment{batch next-state object coordinates}
    \State $p_{ij} \gets \mathrm{softmax}_j\!\left( F_\theta(s_i,a_i,z_i)^\top B_\theta(u_j)/\tau \right)$ per head
    \State $a_i' \gets \pi(s_i', z_i)$, stop-gradient \Comment{one-shot bootstrap action}
    \State $\bar p_{ij}' \gets \mathrm{softmax}_j\!\left( \bar F(s_i',a_i',z_i)^\top \bar B(u_j)/\tau \right)$ per head
    \State $k_{ij} \gets \mathbf 1[j=i]$ (hard) \textbf{or} $\mathrm{softmax}_j\, k_\epsilon(\phi(s_i'), u_j)$ (soft)
    \State $t_{ij} \gets \mathrm{sg}\!\left[ (1-\gamma) k_{ij} + \gamma \bar p_{ij}' \right]$
    \State $\theta \gets \theta - \eta \nabla_\theta \left( -\tfrac1n \sum_{i,j} t_{ij} \log p_{ij} \right)$ \Comment{mean over heads}
    \State $q_i \gets \min_{\text{heads}} \sum_j p_{ij}\big|_{\mathrm{frozen}\ \theta}\; r_{z_i}(u_j)$ \Comment{$r_{z}$ external: goal kernel or reward samples}
    \State update $\pi$: minimize $-\,\mathbb E[q] \,/\, \mathrm{sg}\!\left(\mathbb E|q - \mathbb E q| + \varepsilon\right) + \beta\,\mathcal L_{\mathrm{BC}}(\pi)$
    \State $(\bar F, \bar B) \gets (1-\zeta)(\bar F, \bar B) + \zeta (F_\theta, B_\theta)$
\EndFor
\Statex \textbf{Deployment} (frozen $F, B, \pi$; pool statistics $\mu_B, C_B$):
\State goal $g$: \quad $\hat z \gets \sqrt d\,\operatorname{normalize}(B(g) - \mu_B)$ \Comment{analytic centered point prompt}
\State reward samples $\{(u_j, r_j)\}_{j=1}^{M}$, $M>1$: \quad $\hat z \gets$ Equation~\eqref{eq:centered-dual} \Comment{centered dual encoder}
\State act with $\pi(s, \hat z)$
\end{algorithmic}
\end{algorithm}
}

\end{document}

